% ============================================================================
% SLIDES BD 2026 - TEMA ESCURO (Fundo Preto)
% Normalização & Dependências Funcionais - Centro de Memória
% ============================================================================
\documentclass[10pt,aspectratio=169]{beamer}

% --- Pacotes ---
\usepackage[utf8]{inputenc}
\usepackage[T1]{fontenc}
\usepackage[brazil]{babel}
\usepackage{amsmath,amssymb}
\usepackage{booktabs}
\usepackage{graphicx}
\usepackage{tikz}
\usepackage{hyperref}
\usepackage{array}
\usepackage{multicol}

% --- Tema Escuro ---
\usetheme{Madrid}
\usecolortheme{whale}
\setbeamertemplate{navigation symbols}{}
\setbeamertemplate{footline}[frame number]

% --- Cores do tema escuro ---
\definecolor{primary}{HTML}{1a365d}
\definecolor{secondary}{HTML}{2b6cb0}
\definecolor{accent}{HTML}{d69e2e}
\definecolor{bgdark}{HTML}{0f1419}
\definecolor{codefg}{HTML}{edf2f7}

\setbeamercolor{structure}{fg=accent}
\setbeamercolor{frametitle}{fg=white,bg=primary}
\setbeamercolor{title}{fg=white}
\setbeamercolor{normal text}{fg=white,bg=bgdark}
\setbeamercolor{block title}{fg=white,bg=secondary}
\setbeamercolor{block body}{fg=white,bg=primary!50!black}

% --- Informações ---
\title{Normalização \& Dependências Funcionais}
\subtitle{Teoria Relacional e Decomposição em 1FN, 2FN e 3FN \\(Caso Centro de Memória)}
\author{
  Arthur de Oliveira Lima Potente \\
  Breno Luiz Silva do Carmo \\
  Isaac Salles Gonçalves \\
  Pedro Henrique Rocha de Andrade
}
\institute{Banco de Dados --- 6º Período \\ IFF --- Bom Jesus do Itabapoana, RJ}
\date{02 de Setembro de 2026}

% --- Início ---
\begin{document}

% Slide 1: Capa
\begin{frame}
  \titlepage
\end{frame}

% Slide 2: Sumário
\begin{frame}{Sumário}
  \begin{enumerate}
    \item Conceitos
    \item Cenário 0FN
    \item 1ª FN
    \item 2ª FN
    \item 3ª FN
    \item Conclusões
  \end{enumerate}
\end{frame}

% Slide 3: Conceitos
\begin{frame}{Conceitos}{O que são Dependências Funcionais?}
  \begin{figure}
    \centering
    \textbf{Teoria de DFs e Anomalias}
  \end{figure}
  \vspace{0.3cm}
  \begin{itemize}
    \item \textbf{Definição} ($X \to Y$): X determina unicamente o valor de Y.
    \item \textbf{Determinante vs Determinado}: X determina; Y depende.
    \item \textbf{Anomalias Clássicas}: Inserção, Exclusão e Alteração.
    \item \textbf{Objetivo Central}: Eliminar redundâncias sem perda de dados.
  \end{itemize}
  \vspace{0.2cm}
  \begin{block}{Fonte}
    Elaboração dos Autores (2026)
  \end{block}
\end{frame}

% Slide 4: Cenário 0FN
\begin{frame}{Cenário 0FN}{O Cenário Não-Normalizado}
  \begin{figure}
    \centering
    \textbf{Relação Universal 0FN}
  \end{figure}
  \vspace{0.3cm}
  \begin{itemize}
    \item \textbf{Domínio}: Acervo, doações e exposições unificados.
    \item \textbf{Anomalia de Inserção}: Impossível cadastrar doador sem item do acervo.
    \item \textbf{Anomalia de Exclusão}: Cancelar exposição apaga o item do acervo.
    \item \textbf{Anomalia de Alteração}: Atualização repetida de dados do doador.
  \end{itemize}
  \vspace{0.2cm}
  \begin{block}{Fonte}
    Baseado no Sistema Centro de Memória
  \end{block}
\end{frame}

% Slide 5: 1FN
\begin{frame}{1FN}{Primeira Forma Normal --- Atomicidade}
  \begin{figure}
    \centering
    \textbf{Transição para 1FN e 2FN}
  \end{figure}
  \vspace{0.3cm}
  \begin{itemize}
    \item \textbf{Regra da 1FN}: Atributos estritamente atômicos.
    \item \textbf{Problema (0FN)}: \texttt{palavras\_chave} agrupadas na mesma célula.
    \item \textbf{Decomposição}: Extração para tabela associativa:
    \begin{center}
      \texttt{Item\_PalavraChave(\underline{cod\_tombo}, \underline{palavra\_chave})}
    \end{center}
    \item \textbf{Benefício}: Atomicidade e facilidade de indexação.
  \end{itemize}
  \vspace{0.2cm}
  \begin{block}{Fonte}
    Elaboração dos Autores (2026)
  \end{block}
\end{frame}

% Slide 6: 2FN
\begin{frame}{2FN}{Segunda Forma Normal --- Dependência Total}
  \begin{figure}
    \centering
    \textbf{Eliminação de DFs Parciais}
  \end{figure}
  \vspace{0.3cm}
  \begin{itemize}
    \item \textbf{Regra da 2FN}: Estar na 1FN e sem dependências parciais.
    \item \textbf{Problema}: Título depende só de \texttt{cod\_tombo} em chave composta.
    \item \textbf{Decomposição 2FN}:
    \begin{itemize}
      \item \texttt{Exposicao(\underline{codigo\_exposicao}, titulo, ...)}
      \item \texttt{Item\_Acervo(\underline{cod\_tombo}, titulo, ...)}
      \item \texttt{Item\_Exposicao(\underline{cod\_exp}, \underline{cod\_tombo}, pos)}
    \end{itemize}
    \item \textbf{Garantia}: Atributos dependem da chave primária completa.
  \end{itemize}
  \vspace{0.2cm}
  \begin{block}{Fonte}
    Elaboração dos Autores (2026)
  \end{block}
\end{frame}

% Slide 7: 3FN
\begin{frame}{3FN}{Terceira Forma Normal --- Sem DFs Transitivas}
  \begin{figure}
    \centering
    \textbf{Esquema Final em 3FN}
  \end{figure}
  \vspace{0.3cm}
  \begin{itemize}
    \item \textbf{Regra da 3FN}: Estar na 2FN e sem dependências transitivas.
    \item \textbf{Problema}: \texttt{cod\_tombo} $\to$ \texttt{cod\_doacao} $\to$ \texttt{dados\_doador}.
    \item \textbf{Decomposição 3FN}:
    \begin{itemize}
      \item \texttt{Doacao(\underline{cod\_doacao}, data, termo, id\_doador)}
      \item \texttt{Pessoa(\underline{id\_pessoa}, nome, cpf, email, ...)}
    \end{itemize}
    \item \textbf{Vínculo}: \texttt{Item\_Acervo} mantém apenas FK \texttt{cod\_doacao}.
  \end{itemize}
  \vspace{0.2cm}
  \begin{block}{Fonte}
    Elaboração dos Autores (2026)
  \end{block}
\end{frame}

% Slide 8: Conclusões
\begin{frame}{Conclusões}{Esquema Normalizado e Garantias Formais}
  \begin{figure}
    \centering
    \textbf{Relações Especializadas 3FN}
  \end{figure}
  \vspace{0.3cm}
  \begin{itemize}
    \item \textbf{Lossless Join}: Junção natural recompõe dados sem perdas.
    \item \textbf{Preservação de DFs}: Regras garantidas localmente por constraints.
    \item \textbf{Redundância Nula}: Cada fato registrado em uma única tabela.
    \item \textbf{Integridade}: Consistência assegurada por PKs e FKs.
  \end{itemize}
  \vspace{0.2cm}
  \begin{block}{Fonte}
    Elaboração dos Autores (2026)
  \end{block}
\end{frame}

% Slide 9: Comparação
\begin{frame}{Conclusões}{Comparativo de Esquemas}
  \begin{table}
    \centering
    \small
    \begin{tabular}{>{\raggedright}p{0.35\textwidth} >{\centering}p{0.25\textwidth} >{\centering\arraybackslash}p{0.25\textwidth}}
      \toprule
      \textbf{Critério} & \textbf{0FN (Bruto)} & \textbf{3FN (Normalizado)} \\
      \midrule
      Redundância & Alta & Nula \\
      Anomalias & Críticas & Inexistentes \\
      Integridade & Frágil & Máxima \\
      Update & Lento (múltiplas linhas) & Instantâneo (linha única) \\
      \bottomrule
    \end{tabular}
  \end{table}
\end{frame}

% Slide 10: Referências
\begin{frame}{Referências Bibliográficas}
  \footnotesize
  \begin{enumerate}
    \item ELMASRI, R.; NAVATHE, S. B. \textit{Sistemas de Banco de Dados}. 7. ed. São Paulo: Pearson, 2018.
    \item SILBERSCHATZ, A.; KORTH, H. F.; SUDARSHAN, S. \textit{Sistema de Banco de Dados}. 6. ed. Elsevier, 2012.
    \item CODD, Edgar F. A Relational Model of Data for Large Shared Data Banks. \textit{ACM}, 1970.
    \item DATE, C. J. \textit{Introdução a Sistemas de Bancos de Dados}. 8. ed. Elsevier, 2004.
    \item IBM. \textit{Database Normalization: Guia Prático}. IBM Think Topics, 2024.
  \end{enumerate}
\end{frame}

% Slide 11: Obrigado
\begin{frame}{}
  \centering
  \vspace{1.5cm}
  {\Huge \textbf{Obrigado!}}\\[0.5cm]
  {\large Perguntas \& Discussão}\\[1cm]
  \normalsize
  Contatos:\\
  \texttt{pedroiff0@gmail.com}\\
  \texttt{brenooh7@gmail.com}\\
  \texttt{isaacsalles2005@gmail.com}\\
  \texttt{arthurpotente16@gmail.com}\\[0.5cm]
  \small
  \url{www.phrandrade.com/pt-br/resource/engenharia-de-computação}
\end{frame}

\end{document}
